\clearpage
\section{Application code walkthrough}\label{sec:code}
The excerpts keep the repository's line numbers. Click the source label to open
the exact revision. Start with the POST handler below, then follow the service,
database, and security code as needed.

\subsection{Read an HTTP request and create a Student}
A request to \code{POST /api/students} carries login details and a nested Student
profile. The controller turns that request into a call to the application service.

\par\smallskip
{\small\href{\repo/blob/5d22aed7f6d4ef1e4c0b0c0dabf2962ba4b47608/src/main/java/com/example/fmschool/controller/StudentController.java#L31-L35}{\code{StudentController.java}, lines 31--35}}
\begin{lstlisting}[style=source,language=FMJava,firstnumber=31]
@PostMapping
public ResponseEntity<PeopleApiService.StudentView> create(@Valid @RequestBody CreatePersonRequest request) {
    var created = people.createStudent(request);
    return ResponseEntity.created(URI.create("/api/students/" + created.id())).body(created);
}
\end{lstlisting}


\begin{enumerate}
\item \code{@PostMapping} combines with the class path \code{/api/students}.
The request method and path together select this handler.
\item \code{@RequestBody} binds JSON into \code{CreatePersonRequest}.
\code{@Valid} runs field checks before the method executes. The DTO's nested
\code{@Valid} also checks the profile.
\item The controller delegates to \code{people.createStudent(request)}.
The service creates the account and profile, then builds a response view.
\item \code{ResponseEntity.created} returns 201. Its Location header points to
the new Student profile by ID. The response body is the view object.
\end{enumerate}

\par\smallskip
{\small\href{\repo/blob/5d22aed7f6d4ef1e4c0b0c0dabf2962ba4b47608/src/main/java/com/example/fmschool/dto/CreatePersonRequest.java#L6-L11}{\code{CreatePersonRequest.java}, lines 6--11}}
\begin{lstlisting}[style=source,language=FMJava,firstnumber=6]
public record CreatePersonRequest(
        @NotBlank @Size(min = 3, max = 50)
        @Pattern(regexp = "[a-zA-Z0-9._-]+") String username,
        @NotBlank @Size(min = 8, max = 64) String password,
        @NotNull @Valid ProfileForm profile) {
}
\end{lstlisting}

\par\smallskip
{\small\href{\repo/blob/5d22aed7f6d4ef1e4c0b0c0dabf2962ba4b47608/src/main/java/com/example/fmschool/service/PeopleApiService.java#L48-L52}{\code{PeopleApiService.java}, lines 48--52}}
\begin{lstlisting}[style=source,language=FMJava,firstnumber=48]
public StudentView createStudent(CreatePersonRequest request) {
    Long accountId = createIdentity(request, Role.STUDENT);
    profiles.updateManagedProfile(accountId, request.profile());
    return studentView(students.findByAccountId(accountId).orElseThrow());
}
\end{lstlisting}


\textbf{Which ID is which?}
The account ID identifies the login record. The Student profile ID identifies the
row addressed by \code{/api/students/{id}}. The service looks up the profile using
the account ID it just created, then returns the profile view.

\textbf{Follow the GET path.}
The class's list handler calls \code{PeopleApiService.students()}. The service reads
entities with \code{findAll()} and maps each entity to \code{StudentView}. A view
makes the JSON fields explicit, which keeps the password hash out of the response.

\clearpage
\subsection{Map the Student profile and save an update}
An entity represents a database record. The repository provides operations on that
entity. The service decides which record to load and which fields to update.

\par\smallskip
{\small\href{\repo/blob/5d22aed7f6d4ef1e4c0b0c0dabf2962ba4b47608/src/main/java/com/example/fmschool/entity/Student.java#L7-L27}{\code{Student.java}, lines 7--27}}
\begin{lstlisting}[style=source,language=FMJava,firstnumber=7]
@Entity
@Table(name = "students")
public class Student {

    @Id
    @GeneratedValue(strategy = GenerationType.IDENTITY)
    private Long id;

    @Getter
    @Setter
    @com.fasterxml.jackson.annotation.JsonIgnore
    @OneToOne(fetch = FetchType.LAZY)
    @JoinColumn(name = "account_id", unique = true)
    private Account account;

    @Setter
    @Getter
    private String name;
    @Setter
    @Getter
    private String roll;
\end{lstlisting}

\par\smallskip
{\small\href{\repo/blob/5d22aed7f6d4ef1e4c0b0c0dabf2962ba4b47608/src/main/java/com/example/fmschool/repository/StudentRepository.java#L6-L13}{\code{StudentRepository.java}, lines 6--13}}
\begin{lstlisting}[style=source,language=FMJava,firstnumber=6]
public interface StudentRepository extends JpaRepository<Student, Long> {
    java.util.Optional<Student> findByAccountId(Long accountId);

    @org.springframework.data.jpa.repository.Lock(jakarta.persistence.LockModeType.PESSIMISTIC_WRITE)
    @org.springframework.data.jpa.repository.Query("select s from Student s where s.id = :id")
    java.util.Optional<Student> findLockedById(
            @org.springframework.data.repository.query.Param("id") Long id);
}
\end{lstlisting}

\par\smallskip
{\small\href{\repo/blob/5d22aed7f6d4ef1e4c0b0c0dabf2962ba4b47608/src/main/java/com/example/fmschool/service/StudentService.java#L43-L49}{\code{StudentService.java}, lines 43--49}}
\begin{lstlisting}[style=source,language=FMJava,firstnumber=43]
@Transactional
public void updateStudent(Long id, StudentForm form) {
    Student student = getStudentById(id);
    student.setName(form.getName());
    student.setRoll(form.getRoll());
    studentRepository.save(student);
}
\end{lstlisting}


\begin{itemize}
\item \code{@Id} marks the profile's primary key. \code{IDENTITY} lets PostgreSQL
assign the ID. The unique \code{account_id} links at most one Student profile to
the same account. Lazy loading defers fetching that account until needed.
\item \code{JpaRepository<Student, Long>} declares the entity and ID type.
Spring Data provides \code{findById} and \code{save}. It derives
\code{findByAccountId} from the relationship and property names.
\item \code{findLockedById} uses JPQL, which names the Java entity and its
fields. Its pessimistic write lock holds the row until the transaction ends.
That matters for two actions that could try to enroll the same Student.
\item In \code{updateStudent}, a lookup identifies the record, setters change
the in-memory entity, and the save is part of a transaction. The transaction
commits those changes to the database.
\end{itemize}

\clearpage
\subsection{Sign in and choose a landing page}
Authentication answers who the user is. It loads an account by username and
checks the supplied password against the stored hash. The success handler then
chooses the landing page.

\par\smallskip
{\small\href{\repo/blob/5d22aed7f6d4ef1e4c0b0c0dabf2962ba4b47608/src/main/java/com/example/fmschool/config/SecurityConfig.java#L22-L49}{\code{SecurityConfig.java}, lines 22--49}}
\begin{lstlisting}[style=source,language=FMJava,firstnumber=22]
@Bean
public PasswordEncoder passwordEncoder() {
    return PasswordEncoderFactories.createDelegatingPasswordEncoder();
}

@Bean
public UserDetailsService userDetailsService(AccountRepository accounts) {
    return username -> {
        var account = accounts.findByUsername(username)
                .orElseThrow(() -> new UsernameNotFoundException("Account not found"));
        return User.withUsername(account.getUsername())
                .password(account.getPasswordHash())
                .roles(account.getRole().name())
                .build();
    };
}

@Bean
public AuthenticationSuccessHandler roleLoginSuccessHandler() {
    return (request, response, authentication) -> {
        boolean admin = authentication.getAuthorities().stream()
                .anyMatch(a -> a.getAuthority().equals("ROLE_ADMIN"));
        var handler = new SavedRequestAwareAuthenticationSuccessHandler();
        handler.setAlwaysUseDefaultTargetUrl(true);
        handler.setDefaultTargetUrl(admin ? "/" : "/me");
        handler.onAuthenticationSuccess(request, response, authentication);
    };
}
\end{lstlisting}


\begin{enumerate}
\item The encoder supports stored password hashes with an algorithm identifier.
Spring Security uses it to compare passwords during login.
\item \code{UserDetailsService} loads the username, password hash, and role.
A missing account raises \code{UsernameNotFoundException}.
\item \code{roles(...)} gives Spring Security an authority such as
\code{ROLE_TEACHER}. Rules written as \code{hasRole("TEACHER")} check that authority.
\item The success handler checks for \code{ROLE_ADMIN}. It redirects Admin to
\code{/} and Teacher or Student to \code{/me}. The destination comes from the role
Spring loaded from the account.
\end{enumerate}

\textbf{Where the tests touch it.}
Most MockMvc requests set a test identity directly, so they exercise authorization.
The Teacher-deletion integration test goes further by submitting the Student's
real credentials through the form-login filter after deletion.

\clearpage
\subsection{Choose the security rules for a request}
Spring Security uses the first matching filter chain. The API chain has order 1.
The general web chain has order 2.

\par\smallskip
{\small\href{\repo/blob/5d22aed7f6d4ef1e4c0b0c0dabf2962ba4b47608/src/main/java/com/example/fmschool/config/SecurityConfig.java#L51-L86}{\code{SecurityConfig.java}, lines 51--86}}
\begin{lstlisting}[style=source,language=FMJava,firstnumber=51]
@Bean
@Order(2)
public SecurityFilterChain securityFilterChain(HttpSecurity http) throws Exception {
    return http
            .authorizeHttpRequests(auth -> auth
                    .dispatcherTypeMatchers(DispatcherType.ERROR, DispatcherType.FORWARD).permitAll()
                    .requestMatchers("/login", "/signup", "/signup/student", "/signup/teacher",
                            "/css/**", "/favicon.ico").permitAll()
                    .requestMatchers("/students", "/students/**",
                            "/api/students", "/api/students/**",
                            "/admin", "/admin/**").hasRole("ADMIN")
                    .anyRequest().authenticated())
            .formLogin(form -> form.loginPage("/login")
                    .successHandler(roleLoginSuccessHandler()).permitAll())
            .httpBasic(Customizer.withDefaults())
            .logout(logout -> logout.logoutSuccessUrl("/login?logout"))
            .build();
}

@Bean
@Order(1)
public SecurityFilterChain apiFilterChain(HttpSecurity http) throws Exception {
    return http.securityMatcher("/api/**")
            .authorizeHttpRequests(auth -> auth
                    .requestMatchers("/api/students", "/api/students/**",
                            "/api/teachers", "/api/teachers/**").hasRole("ADMIN")
                    .anyRequest().authenticated())
            .httpBasic(basic -> basic.authenticationEntryPoint((request, response, exception) ->
                    apiError(response, 401, "Authentication required")))
            .exceptionHandling(errors -> errors
                    .authenticationEntryPoint((request, response, exception) ->
                            apiError(response, 401, "Authentication required"))
                    .accessDeniedHandler((request, response, exception) ->
                            apiError(response, 403, "Access denied or missing CSRF token")))
            .build();
}
\end{lstlisting}


\begin{itemize}
\item \code{securityMatcher("/api/**")} selects the API chain. The route rules
inside it restrict the people collections to Admin and require authentication for
the other API routes. Basic authentication is available for API clients.
\item The web chain opens login, signup, and static CSS. It protects Admin routes
and configures the form login and logout pages.
\item \code{401} means authentication is required. \code{403} can mean a role or
ownership check failed, or a write lacked a valid CSRF token.
\end{itemize}

A route can permit a Teacher to call an enrollment endpoint. The service must still
check whether the specific enrollment belongs to that Teacher. The next two
sections show that distinction in code.

\clearpage
\subsection{Apply, invite, and respond to an enrollment}
A Student application and a Teacher invitation use the same enrollment table.
\code{requestedBy} records which side sent the request, so the service can identify
the recipient. A rejected pair can reopen through the shared \code{request} method.

\par\smallskip
{\small\href{\repo/blob/5d22aed7f6d4ef1e4c0b0c0dabf2962ba4b47608/src/main/java/com/example/fmschool/service/EnrollmentService.java#L105-L121}{\code{EnrollmentService.java}, lines 105--121}}
\begin{lstlisting}[style=source,language=FMJava,firstnumber=105]
@Transactional
public void respond(Long id, boolean accept) {
    Account actor = currentAccount();
    Enrollment enrollment = enrollment(id);
    // Check who may act before exposing the current request state.
    if (actor.getRole() != Role.ADMIN && !isRecipient(enrollment, actor)) {
        throw forbidden();
    }
    if (enrollment.getStatus() != EnrollmentStatus.PENDING) {
        throw conflict("Only pending requests can be accepted or rejected");
    }
    if (accept) {
        requireAvailable(enrollment.getStudent().getId(), enrollment.getTeacher().getId());
    }
    enrollment.setStatus(accept ? EnrollmentStatus.ACCEPTED : EnrollmentStatus.REJECTED);
    enrollments.save(enrollment);
}
\end{lstlisting}

\par\smallskip
{\small\href{\repo/blob/5d22aed7f6d4ef1e4c0b0c0dabf2962ba4b47608/src/main/java/com/example/fmschool/service/EnrollmentService.java#L136-L153}{\code{EnrollmentService.java}, lines 136--153}}
\begin{lstlisting}[style=source,language=FMJava,firstnumber=136]
private void request(Teacher teacher, Student student, Role sender) {
    student = lockStudent(student.getId());
    requireAvailable(student.getId(), teacher.getId());
    Enrollment enrollment = enrollments.findLockedPair(teacher.getId(), student.getId())
            .orElseGet(Enrollment::new);
    if (enrollment.getStatus() == EnrollmentStatus.ACCEPTED) {
        throw conflict("This Student is already associated with this Teacher");
    }
    if (enrollment.getStatus() == EnrollmentStatus.PENDING) {
        throw conflict("A request is already pending; its recipient must respond");
    }
    // A rejected pair can apply/invite again, without duplicating the association.
    enrollment.setTeacher(teacher);
    enrollment.setStudent(student);
    enrollment.setRequestedBy(sender);
    enrollment.setStatus(EnrollmentStatus.PENDING);
    enrollments.save(enrollment);
}
\end{lstlisting}


\begin{itemize}
\item \code{respond} checks the actor before checking the request status. Only Admin
or the recipient can respond. A second acceptance encounters the status guard
and returns a conflict before the save.
\item Accepting also calls \code{requireAvailable}. It checks whether the Student
already has an accepted enrollment with another Teacher.
\item \code{request} locks the Student and finds the existing Teacher--Student pair.
For a rejected request it reuses the row, sets the new sender, and returns the
status to pending. A pending pair already has a recipient who needs to respond.
\end{itemize}

\textbf{Concurrency.}
\code{enrollment(id)} locks the Student before the enrollment. That serializes
competing acceptances for the same Student. The database also enforces the
one-accepted-course rule through a partial unique index on \code{student_id} for
accepted rows. The course transfer test exercises the rule sequentially.

\clearpage
\subsection{Edit marks through the accepted enrollment}
Teacher edit access depends on the role, the enrollment's status, and the
account that owns the Teacher profile. The service uses the enrollment ID to
load that relationship before writing either record.

\par\smallskip
{\small\href{\repo/blob/5d22aed7f6d4ef1e4c0b0c0dabf2962ba4b47608/src/main/java/com/example/fmschool/service/EnrollmentService.java#L204-L225}{\code{EnrollmentService.java}, lines 204--225}}
\begin{lstlisting}[style=source,language=FMJava,firstnumber=204]
@Transactional
@PreAuthorize("hasAnyRole('TEACHER', 'ADMIN')")
public void updateEnrolledStudent(Long id, @Valid EnrolledStudentForm form) {
    Enrollment enrollment = enrollment(id);
    requireEditor(enrollment, currentAccount());
    Student student = enrollment.getStudent();
    student.setName(form.getName());
    student.setRoll(form.getRoll());
    enrollment.setMarks(form.getMarks());
    students.save(student);
    enrollments.save(enrollment);
    // Account username/password/role are deliberately untouched.
}

private void requireEditor(Enrollment enrollment, Account actor) {
    if (enrollment.getStatus() != EnrollmentStatus.ACCEPTED
            || (actor.getRole() != Role.ADMIN
                && !(actor.getRole() == Role.TEACHER
                    && Objects.equals(enrollment.getTeacher().getAccount().getId(), actor.getId())))) {
        throw forbidden();
    }
}
\end{lstlisting}

\par\smallskip
{\small\href{\repo/blob/5d22aed7f6d4ef1e4c0b0c0dabf2962ba4b47608/src/main/resources/db/changelog/db.changelog-master.yaml#L170-L173}{\code{db.changelog-master.yaml}, lines 170--173}}
\begin{lstlisting}[style=source,language=YAML,firstnumber=170]
    sql: CREATE UNIQUE INDEX enrollments_one_accepted_course ON enrollments (student_id) WHERE status = 'ACCEPTED'
- sql:
    dbms: postgresql
    sql: ALTER TABLE enrollments ADD CONSTRAINT enrollments_marks_range CHECK (marks IS NULL OR (marks >= 0 AND marks <= 100))
\end{lstlisting}


\begin{enumerate}
\item \code{@PreAuthorize} permits Teacher and Admin roles through the method
boundary. \code{requireEditor} requires an accepted enrollment. For Teachers
it compares the enrollment's Teacher account ID with the actor's account ID.
\item The name and roll belong to the Student profile. The mark belongs to the
enrollment, so both are updated within the transaction. Login details stay in
the linked account record.
\item \code{form.getMarks()} can return null. The service stores that null as
ungraded. Zero remains a valid numeric mark. The schema permits null or a
value between 0 and 100.
\end{enumerate}

\textbf{What the board shows.}
\code{getBoard()} gets enrollment rows separately from the lists of people the
user can invite or apply to. That distinction lets a Student see pending history
even when an accepted enrollment hides the new-application picker. A Teacher's
picker filters out Students who already have an accepted course or have no account.

\clearpage
\subsection{Roll back creation and clean up deletions}
The Admin API works across two kinds of records. Account creation writes the
identity first and profile creation fills in the role-specific record afterward.

\par\smallskip
{\small\href{\repo/blob/5d22aed7f6d4ef1e4c0b0c0dabf2962ba4b47608/src/main/java/com/example/fmschool/service/PeopleApiService.java#L13-L16}{\code{PeopleApiService.java}, lines 13--16}}
\begin{lstlisting}[style=source,language=FMJava,firstnumber=13]
@Service
@PreAuthorize("hasRole('ADMIN')")
@Transactional
public class PeopleApiService {
\end{lstlisting}

\par\smallskip
{\small\href{\repo/blob/5d22aed7f6d4ef1e4c0b0c0dabf2962ba4b47608/src/main/java/com/example/fmschool/service/PeopleApiService.java#L78-L82}{\code{PeopleApiService.java}, lines 78--82}}
\begin{lstlisting}[style=source,language=FMJava,firstnumber=78]
public TeacherView createTeacher(CreatePersonRequest request) {
    Long accountId = createIdentity(request, Role.TEACHER);
    profiles.updateManagedProfile(accountId, request.profile());
    return teacherView(teachers.findByAccountId(accountId).orElseThrow());
}
\end{lstlisting}

\par\smallskip
{\small\href{\repo/blob/5d22aed7f6d4ef1e4c0b0c0dabf2962ba4b47608/src/main/java/com/example/fmschool/service/ProfileService.java#L142-L150}{\code{ProfileService.java}, lines 142--150}}
\begin{lstlisting}[style=source,language=FMJava,firstnumber=142]
} else if (account.getRole() == Role.TEACHER) {
    if (form.getCourse() == null || form.getCourse().isBlank()) {
        throw new IllegalArgumentException("Course is required");
    }
    Teacher teacher = teachers.findByAccountId(account.getId()).orElseGet(Teacher::new);
    teacher.setAccount(account);
    teacher.setName(form.getName());
    teacher.setCourse(form.getCourse());
    teachers.save(teacher);
\end{lstlisting}

\par\smallskip
{\small\href{\repo/blob/5d22aed7f6d4ef1e4c0b0c0dabf2962ba4b47608/src/main/java/com/example/fmschool/service/ProfileService.java#L156-L163}{\code{ProfileService.java}, lines 156--163}}
\begin{lstlisting}[style=source,language=FMJava,firstnumber=156]
private void delete(Account account) {
    students.findByAccountId(account.getId()).ifPresent(students::delete);
    teachers.findByAccountId(account.getId()).ifPresent(teachers::delete);
    // Remove dependent records before the account, preserving unrelated people.
    students.flush();
    teachers.flush();
    accounts.delete(account);
}
\end{lstlisting}


\begin{itemize}
\item The class-level \code{@Transactional} on \code{PeopleApiService} gives
\code{createTeacher} a transaction that covers both steps. The nested account
and profile service calls join that transaction.
\item A missing course reaches the Teacher branch of \code{save} and throws an
\code{IllegalArgumentException}. The API handler returns 400. The transaction
rolls back, so the account written by the first step disappears.
\item Deleting an account removes its profiles first and flushes those deletions.
The enrollment foreign keys have \code{ON DELETE CASCADE}, so the linked
enrollment rows disappear when the Teacher profile is deleted.
\item The Student profile belongs to another account. The cascade follows the
deleted Teacher's enrollment links and leaves the Student's login usable.
\end{itemize}

\textbf{What to show in the lab.}
Open integration test I2 for the rollback query and I3 for the row checks plus the
Student's login after the Teacher is gone.

\clearpage
\subsection{Connect the web form and the database}
Thymeleaf renders the profile form with values supplied by Spring MVC. The
browser submits the form to the route shown in the generated HTML.

\par\smallskip
{\small\href{\repo/blob/5d22aed7f6d4ef1e4c0b0c0dabf2962ba4b47608/src/main/resources/templates/profile.html#L27-L51}{\code{profile.html}, lines 27--51}}
\begin{lstlisting}[style=source,language=HTML,firstnumber=27]
<h2>Personal information</h2>
<form
  th:action="@{${savePath}}"
  th:object="${profileForm}"
  method="post"
  class="field-form"
>
  <ul role="alert" th:if="${#fields.hasAnyErrors()}">
    <li
      th:each="error : ${#fields.allErrors()}"
      th:text="${error}"
    ></li>
  </ul>
  <p>
    <label for="name">Name</label>
    <input id="name" type="text" th:field="*{name}" />
  </p>
  <p th:if="${isStudent}">
    <label for="roll">Roll</label>
    <input
      id="roll"
      inputmode="numeric"
      type="text"
      th:field="*{roll}"
    />
\end{lstlisting}


\begin{itemize}
\item \code{th:action} generates the route from the controller's save path.
\code{th:object} selects the profile form object. \code{th:field} connects an
input to one of its fields.
\item The error list displays validation messages after invalid input.
\code{isStudent} selects the Student roll input. The server also validates the
profile before saving it.
\item The shared layout loads the compiled Tailwind stylesheet from
\code{/css/app.css}. Changing Tailwind source requires a CSS build.
\end{itemize}

\par\smallskip
{\small\href{\repo/blob/5d22aed7f6d4ef1e4c0b0c0dabf2962ba4b47608/compose.yaml#L1-L14}{\code{compose.yaml}, lines 1--14}}
\begin{lstlisting}[style=source,language=YAML,firstnumber=1]
services:
  postgres:
    image: 'postgres:18.6'
    environment:
      - 'POSTGRES_DB=fm_school_db'
      - 'POSTGRES_PASSWORD=secret'
      - 'POSTGRES_USER=myuser'
    ports:
      - '127.0.0.1:5432:5432'
    volumes:
      - 'postgres_data:/var/lib/postgresql'

volumes:
  postgres_data:
\end{lstlisting}


The PostgreSQL service publishes its port to localhost. Spring Boot runs on the
host and uses its PostgreSQL driver to reach that port. The named volume keeps
database files across ordinary container recreation.

\clearpage
\subsection{Read the GitHub Actions workflow}
The workflow runs the same Maven verification goal used in local development.
GitHub starts PostgreSQL in a service container on the Ubuntu runner.

\par\smallskip
{\small\href{\repo/blob/5d22aed7f6d4ef1e4c0b0c0dabf2962ba4b47608/.github/workflows/tests.yml#L1-L38}{\code{tests.yml}, lines 1--38}}
\begin{lstlisting}[style=source,language=YAML,firstnumber=1]
name: FM School tests

on:
  push:
  pull_request:
  workflow_dispatch:

permissions:
  contents: read

jobs:
  verify:
    runs-on: ubuntu-latest
    timeout-minutes: 12
    services:
      postgres:
        image: postgres:18.6
        env:
          POSTGRES_DB: fm_school_test
          POSTGRES_USER: fm_school_test
          POSTGRES_PASSWORD: fm_school_test
        ports:
          - 55432:5432
        options: >-
          --health-cmd "pg_isready -U fm_school_test -d fm_school_test"
          --health-interval 5s
          --health-timeout 5s
          --health-retries 12

    steps:
      - uses: actions/checkout@v6
      - uses: actions/setup-java@v5
        with:
          distribution: temurin
          java-version: '25'
          cache: maven
      - name: Run unit and PostgreSQL integration tests
        run: mvn --batch-mode --no-transfer-progress verify
\end{lstlisting}


\begin{enumerate}
\item \code{push} and \code{pull_request} schedule the checks. The workflow also
supports a manual run.
\item The PostgreSQL service waits for \code{pg_isready}. Port 55432 on the runner
maps to port 5432 inside the test database container.
\item Checkout retrieves the repository. Setup Java installs Temurin 25 and
caches Maven dependencies.
\item \code{mvn verify} runs the isolated tests, packages the application,
runs the PostgreSQL tests, and fails the job if either phase reports a failure.
\end{enumerate}

The integration profile uses the same database name and port as this service.
The workflow's successful \code{verify} job is the result to show during the demo.
